\documentclass[10pt,a4paper]{article}
\usepackage{amsmath,amssymb,amsthm,graphicx}
\newcommand{\rk}{\operatorname{rk}}
\newcommand{\xc}{\operatorname{xc}}
\newcommand{\sep}{\operatorname{sep}}
\newcommand{\SGN}{\operatorname{SGN}}
\newcommand{\sgn}{\operatorname{sgn}}
\newcommand{\conv}{\operatorname{conv}}
\newcommand{\RR}{\mathbb R}
\newenvironment{originalresult}[1]{\par\medskip\noindent\textbf{#1.}\itshape}{\par\medskip}
\title{Hard submatrices for non-negative rank}
\author{Pavel Hrubeš}
\date{doi:10.4230/LIPIcs.CCC.2024.13; CC BY 4.0}
\begin{document}
\maketitle
\noindent\textit{Editorial note.} Faithful transcription from official CCC 2024 CC BY 4.0 PDF; original notation, indexing and dimensions retained. Author TeX was not located. No author proof is replaced or extended; standard cited prior results are retained as prerequisites, and bibliography entries keep their original numbers.\par\medskip
\noindent\textit{Allowed standard prerequisites (editorial).}\begin{itemize}
\item Slack-matrix/extension-complexity sandwich bound. rk\_+(S)-1 <= xc(P0,P1) <= rk\_+(S). Braun et al. [2], source Lemma 14; exact statement retained
\item Single existential-block quantifier elimination bound. Basu, Pollack and Roy [1], Theorem 14.16; exact used statement retained in appendix
\end{itemize}\medskip
\section*{3 A combinatorial lemma (author source pp. 13:3--13:4)}
Both Theorems 1 and 2 rely on a simple combinatorial lemma.
\begin{originalresult}{Lemma 3}
Let $\mathcal A\subseteq 2^{[n]}$ be a family of subsets of $[n]$. Assume that $1\le k\le n$ is such that every $k$-element subset of $[n]$ is contained in some $A\in\mathcal A$. Then there exists a subfamily $\mathcal A'\subseteq\mathcal A$ of size $|\mathcal A'|\le O(|\mathcal A|^{1/k}\log(n/k))$ with $\bigcup\mathcal A'=[n]$. In particular, if $|\mathcal A|\le 2^k$ then $|\mathcal A'|\le O(\log n)$.
\end{originalresult}
\begin{proof}
Assume that $|\mathcal A|\le a^k$. Let $t$ be the size of a largest set in $\mathcal A$. Then we have
\[\binom nk\le a^k\binom tk.\]
Hence $t\ge\frac n{ea}$, using the estimates $(\frac nk)^k\le\binom nk$, $\binom tk\le(\frac{et}k)^k$. Take some $A_0\in\mathcal A$ of size $t$. Let
\[\mathcal A_1:=\{A\setminus A_0:A\in\mathcal A\}.\]
Then every subset of $U_1:=[n]\setminus A_0$ of size at most $k$ is contained in some member of $\mathcal A_1$. The size of $U_1$ is at most $n(1-\frac1{ea})$. Similarly, take a largest set $A_1$ from $\mathcal A_1$ and obtain a new family $\mathcal A_2\subseteq 2^{U_2}$ on $U_2:=U_1\setminus A_1$. After $s$ steps, the size of $U_s$ is at most $n(1-\frac1{ea})^s$ and after $s\le O(a\log(n/k))$ steps we have $|U_s|\le k$. This guarantees that the largest set in $\mathcal A_s$ is $U_s$ itself and $[n]=\bigcup_{i=0}^s A_s$. By construction, every $A_i$ is contained in some element of the original family $\mathcal A$.
\end{proof}

\section*{6.1 Extended formulations and separation complexity (author source p. 13:8)}
Let us first make a short detour into extended formulations of convex polyhedra.
A polyhedron $P\subseteq\RR^r$ is a (possibly unbounded) set defined by a finite number of linear constraints. Following [20,18,2], define the extension complexity of $P$, $\xc(P)$, as the smallest $s$ such that $P$ is a linear projection of a polyhedron $Q\subseteq\RR^m$ where $Q$ can be defined using $s$ inequalities (and any number of equalities). Observe that $P$ with extension complexity $s$ can be expressed in the standard form
\[x\in P\text{ iff }\exists y\in\RR^s\ Cx+Dy=b,\ y\ge0,\]
where $C\in\RR^{t\times r}$, $D\in\RR^{t\times s}$ and $b\in\RR^t$ for some $t$.
Let $V$ be a finite subset of $\RR^r$. Given $A\subseteq V$, its separation complexity, $\sep_V(A)$, is the minimum $\xc(P)$ over all polyhedra $P\subseteq\RR^r$ with\footnote{If no such polyhedron exists, which may happen if $V$ is not convexly independent, we set $\sep_V(A):=\infty$.}
\[P\cap V=A;\]
such a $P$ is called a separating polyhedron for $A$. In other words, $\sep_V(A)$ is the smallest $s$ so that we can distinguish points in $A$ from points in $V\setminus A$ by means of a linear program with $s$ inequalities. Moreover, such a program can be rewritten as
\[x\in A\text{ iff }(x\in V\text{ and }\exists y\in\RR^s\ Cx+Dy=b,\ y\ge0).\]
The notion of separation complexity has been studied in [7,8,9] in the case when $V=\{0,1\}^n$ is the Boolean cube. The following theorem is of independent interest and can be seen as an extension of similar results in [7,9]. The proof is a considerable simplification of the previous ones.

\begin{originalresult}{Theorem 13 (author source p. 13:8)}
Let $V$ be a non-empty finite subset of $\RR^r$. Given a parameter $s\ge1$, let $\mathcal A$ be the collection of subsets $A$ of $V$ with $\sep_V(A)\le s$. Then
\[|\mathcal A|\le2^{O(s(r+s)^2\log|V|)}.\]
\end{originalresult}

\section*{A Proof of Theorem 13 (author source pp. 13:11--13:12)}
The proof uses known results on quantifier elimination which we first outline. We follow the monograph of Basu, Pollack and Roy [1]. Theorem 13 requires an elimination of only a single block of existential quantifiers, so we focus on this case only.
For $b\in\RR$, let
\[\sgn(b):=\begin{cases}1,&b>0,\\0,&b=0,\\-1,&b<0.\end{cases}\]
Given $b=\langle b_1,\ldots,b_m\rangle\in\RR^m$, let $\sgn(b):=\langle\sgn(b_1),\ldots,\sgn(b_m)\rangle\in\{-1,0,1\}^m$. Let $F=F(z,y)$ be a sequence of $m$ polynomials $f_1,\ldots,f_m\in\RR[z,y]$ in variables $z=\{z_1,\ldots,z_{k_1}\}$ and $y=\{y_1,\ldots,y_{k_2}\}$. Given $a\in\RR^{k_1}$, define $\SGN_1(F,a)\subseteq\{-1,0,1\}^m$
\[\SGN_1(F,a):=\{\sgn(F(a,b)):b\in\RR^{k_2}\}.\]
Let
\[\SGN(F):=\{\SGN_1(F,a):a\in\RR^{k_1}\}.\]
Theorem 14.16 from [1] provides the following bound on the size of SGN:
\begin{originalresult}{Theorem ([1])}
If every polynomial in $F$ has degree at most $d$ then
\[|\SGN(F)|\le m^{(k_1+1)(k_2+1)}d^{O(k_1)O(k_2)}.\tag{4}\]
\end{originalresult}
We now apply this result to the case of Theorem 13. Let $V,s,\mathcal A$ be as in the assumption. Every $A\in\mathcal A$ can be described by a linear system with $s$ inequalities. Namely, for every $x\in V$,
\[x\in A\text{ iff }\exists y\in\RR^s\ Cx+Dy=b,\ y\ge0,\tag{5}\]
where $C\in\RR^{t\times r}$, $D\in\RR^{t\times s}$ and $b\in\RR^t$. Since $Cx+Dy=b$ is a system of equations in $r+s$ variables $x,y$, we can also assume $t=r+s$.
Let us view the parameters $C,D,b$ in (5) as variables. Let $z$ be the set of these variables, of size $k_1=(r+s)(r+s+1)$. Given $v\in V$, let $F_v(z,y)$ be the sequence of $(r+s)$ polynomials
\[Cv+Dy-b\]
in variables $z$ and $y=\{y_1,\ldots,y_s\}$. Let $F(z,y)$ be the union of $F_v(z,y)$ over all $v\in V$, together with the polynomials $y_1,\ldots,y_s$. Hence $F$ consists of $m=s+|V|(r+s)$ polynomials of degree at most two.
$F(z,y)$ is set up so that
\[|\mathcal A|\le|\SGN(F)|.\]
To see this, observe that whenever the parameters $z$ are fixed, the set $A\subseteq V$ given by (5) is uniquely determined by $\SGN_1(F(z,y))$. Since every $A\in\mathcal A$ is obtained by some fixing of the parameters, we indeed obtain $|\mathcal A|\le|\SGN(F(z,y))|$.
Finally, we can apply (4) to estimate $|\SGN(F)|$ with $m=s+|V|(r+s)$, $k_1=(r+s)(r+s+1)$, $k_2=s$, and $d=2$. To simplify the expression, we can assume $s+r\le|V|$; otherwise the upper bound asserted in Theorem 13 exceeds the trivial bound $|\mathcal A|\le2^{|V|}$. This means that $m\le2|V|^2$. If we loosen the bound (4) as $|\SGN(F)|\le(dm)^{O(k_1)O(k_2)}$, we obtain (recall that $s\ge1$)
\[|\SGN(F)|\le2^{O(s(s+r)^2\log|V|)},\]
as required.

\section*{Non-negative rank (author source p. 13:2)}
Given an $m\times n$ matrix $M$ with real non-negative entries, its non-negative rank, $\rk_+(M)$, is the smallest $s$ such that $M$ can be written as
\[M=LR,\]
where $L$ and $R$ are non-negative matrices of dimensions $m\times s$ and $s\times n$, respectively.
\paragraph{Notation.} All logarithms are in base 2 and $[n]:=\{1,\ldots,n\}$.
\section*{6.2 Submatrices of large non-negative rank (author source pp. 13:8--13:9)}
In order to apply Theorem 13, we also need a connection between extension complexity and non-negative rank. This is provided by the notion of slack matrix introduced in [20]. Following [20,2], we now define what it is. Let $V$ be a sequence $v_1,\ldots,v_{m_1}$ of points in $\RR^r$ and $L(x)$ a system $\ell_1(x)\ge b_1,\ldots,\ell_{m_2}(x)\ge b_{m_2}$ of inequalities in $\RR^r$. The slack matrix with respect to $V$ and $L(x)$ is the $m_2\times m_1$ matrix $S$ such that
\[S_{i,j}=\ell_i(v_j)-b_i.\]
Let $P_0:=\conv(V)$ be the convex hull of $V$ and $P_1:=\{x\in\RR^r:L(x)\text{ holds}\}$. If $P_0\subseteq P_1$ then $S$ is non-negative. In [2], we can find:
\begin{originalresult}{Lemma 14 ([2])}
Let $P_0\subseteq P_1$ and $S$ be as above. Define $\xc(P_0,P_1)$ as the minimum $\xc(P)$ over all polyhedra with $P_0\subseteq P\subseteq P_1$. Then
\[\rk_+S-1\le\xc(P_0,P_1)\le\rk_+S.\]
\end{originalresult}

\paragraph{Subadditivity (author source p.13:4).}
Observe that\footnote{If $A_1,A_2$ are disjoint, this is quite obvious. Otherwise consider $A_1,A_2\setminus A_1$.}
\[\chi_1(M_{A_1\cup A_2})\le\chi_1(M_{A_1})+\chi_1(M_{A_2}),\tag{1}\]
and so $\chi_1(M_A)$ can be viewed as a subadditive measure on $[n]$ whenever $M$ is fixed.

\begin{originalresult}{Theorem 1 (restated; author source p. 13:9)}
Let $M$ be an $m\times n$ non-negative real matrix with $n\ge2$. Then for every $k\le n$, $M$ contains an $m\times k$ submatrix of $k$ columns with non-negative rank $\Omega(R)$, where $R:=\min((\frac k{\log n})^{1/3},\frac{\rk_+(M)}{\log n})$.
\end{originalresult}

\begin{proof}
Let $r$ be the rank of $M$. We can write $M=LR$ where $L\in\RR^{m\times r}$, $R\in\RR^{r\times n}$. Let $V\subseteq\RR^r$ be the set of columns $v_1,\ldots,v_n$ of $R$. (Without loss of generality, the columns of $M$ are distinct). Given $A\subseteq[n]$, let $M_A$ be the submatrix obtained by deleting columns outside of $A$ from $M$. Also let $V_A:=\{v_i:i\in A\}$. Then $M_A$ can be interpreted as the slack matrix of the polytope $P_A=\conv(V_A)$ and the polyhedron $Q=\{x\in\RR^d:Lx\ge0\}$.
Suppose that for every $A$ of size $k$, $\rk_+(M_A)\le s$. Then for every such $A$, there is a polyhedron $Q_A$ with $V_A\subseteq Q_A\subseteq Q$ with $\xc(Q_A)\le s$. Let $A^*:=V\cap Q_A$. Then $Q_A$ is a separating polyhedron for $A^*\supseteq A$. Let $\mathcal A$ be the collection of $A^*$ over all $A$ of size $k$. Theorem 13 implies
\[|\mathcal A|\le2^{c\log n(s+r)^3},\]
where $c$ is an absolute constant.
We will assume $r\le(\frac k{2c\log n})^{1/3}$. Otherwise $M$ contains a full rank $r\times r$ submatrix, $\rk_+$ is lower-bounded by rank, and the conclusion of the theorem follows.
If $|\mathcal A|\ge2^k$, we obtain $c\log n(s+r)^3\ge k$ and hence $s\ge\Omega((k/\log n)^{1/3})$ from the assumption on $r$.
Assume $|\mathcal A|\le2^k$. By Lemma 3, there exists a subfamily $\mathcal A'\subseteq\mathcal A$ of size $O(\log n)$ which covers $[n]$. This implies (note that (1) holds also for non-negative rank) $\rk_+(M)\le O(s\log n)$ and $s\ge\Omega(\rk_+(M)/\log n)$.
\end{proof}

\begin{thebibliography}{99}

\bibitem[1]{source1} S. Basu, R. Pollack, and M.F. Roy. Algorithms in real algebraic geometry. Springer-Verlag, 2006.
\bibitem[2]{source2} G. Braun, S. Fiorini, S. Pokutta, and D. Steuer. Approximation limits of linear programs (beyond hierarchies). Math. of Operations Research, 40(3), 2015.
\bibitem[3]{source3} S. Chandran, D. Issac, and A. Karrenbauer. On the parameterized complexity of biclique cover and partition. In International Symposium on Parameterized and Exact Computation, 2016.
\bibitem[7]{source7} P. Hrube\v{s}. On $\epsilon$-sensitive monotone computations. Computational Complexity, 29(2), 2020.
\bibitem[8]{source8} P. Hrube\v{s}. On the complexity of computing a random Boolean function over the reals. Theory of Computing, 16(9):1--12, 2020.
\bibitem[9]{source9} P. Hrube\v{s} and N. Talebanfard. On the extension complexity of polytopes separating subsets of the Boolean cube. Disc. and Comp. Geom., 70(1), 2022.
\bibitem[18]{source18} Thomas Rothvo\ss. Some 0/1 polytopes need exponential size extended formulations. CoRR, abs/1105.0036, 2011. arXiv:1105.0036.
\bibitem[20]{source20} Mihalis Yannakakis. Expressing combinatorial optimization problems by linear programs. Journal of Computer and System Sciences, 43(3):441--466, 1991.

\end{thebibliography}
\end{document}
